\begin{table}[t]
\centering
\caption{Summary Statistics}
\label{tab:summary}
\begin{tabular}{lcc}
\toprule
 & White & Black \\
\midrule
Occ. Score 1920 & 15.54 & 11.04 \\
Occ. Score 1930 & 24.39 & 17.02 \\
Occ. Score 1940 & 25.37 & 17.99 \\
$\Delta$ Occ. 1930--40 & 0.98 & 0.97 \\
SEI 1930 & 29.66 & 14.29 \\
SEI 1940 & 30.15 & 14.72 \\
Age 1930 & 34.6 & 32.9 \\
Farm Worker & 0.282 & 0.464 \\
ND Spending p.c. & 132.9 & 101.4 \\
N & 10,641,469 & 394,945 \\
\bottomrule
\end{tabular}
\begin{tablenotes}
\item \textit{Notes:} Sample consists of men aged 18--55 in 1930, linked across the 1920, 1930, and 1940 full-count censuses via the IPUMS Machine Learning Panel (MLP v2). Occupational score (OCCSCORE) is the Duncan socioeconomic index mapping 1950-coded occupations to median income. New Deal spending per capita is total federal grants (WPA, FERA, CWA, and other programs) from Fishback, Kantor, and Wallis (2003) at the county level. Farm Worker = share in agricultural occupations in 1930.
\end{tablenotes}
\end{table}
